\documentclass[11pt,a4paper]{article}
\usepackage[utf8]{inputenc}
\usepackage[T1]{fontenc}
\usepackage{amsmath,amssymb}
\usepackage[margin=2.5cm]{geometry}
\usepackage{xcolor}

\setlength{\parindent}{0pt}
\setlength{\parskip}{6pt}

\begin{document}

\begin{center}
{\large Technical Neural Networks Assignment Sheet 6}\\[4pt]
{\bfseries Names: Mohamadhosein Sabaghi (50547423)}
\end{center}

{\small (Programming Assignment PA-C is solved separately.)}

\section*{Assignment 27}

\textbf{Common aspects.}
\begin{itemize}
  \item Both are feed-forward networks trained in a supervised way.
  \item Both are universal function approximators, usable for regression and classification.
  \item Both compute their output from simple processing units connected by weights, organized
  in layers (input $\to$ hidden/RBF $\to$ output).
\end{itemize}

\textbf{Different aspects.}
\begin{itemize}
  \item \textbf{Depth/structure:} RBF networks always have exactly 3 layers; MLPs can have 2 or
  more layers.
  \item \textbf{Hidden-unit response:} MLP hidden neurons apply a transfer function to a
  \emph{weighted sum} of inputs -- their response extends along a hyperplane, i.e.\ it is
  \emph{global}. RBF neurons apply a unimodal function to the \emph{distance to a center} -- their
  response is \emph{local}, decaying to $\approx 0$ away from their own center.
  \item \textbf{Output layer:} RBF output neurons are always linear (identity transfer function,
  just a weighted sum); MLP layers (including usually the output) are typically nonlinear.
  \item \textbf{Training:} RBF training splits into two simpler, largely independent
  sub-problems -- setting the centers/widths, then solving for the output weights, which is a
  \emph{linear} least-squares problem (exactly solvable, e.g.\ via the Moore-Penrose
  pseudo-inverse, assignment 29). MLP training requires nonlinear Backpropagation of Error through
  every layer at once -- a nonconvex problem with local minima and more hyperparameter tuning.
  \item \textbf{Interpretability:} RBF centers are interpretable as prototype points in input
  space; MLP hidden units (hyperplanes) are less directly interpretable.
  \item \textbf{Scaling with input dimension:} covering an $N$-dimensional space with local
  receptive fields needs a number of RBF neurons that grows rapidly with $N$ (curse of
  dimensionality hits RBF harder); MLP hyperplane-based hidden units scale better to high $N$.
  \item \textbf{Extrapolation:} RBF outputs fall back to $\approx 0$ (or BIAS) far from all
  centers -- safe but uninformative outside the training range. MLP units don't vanish away from
  training data, so they can extrapolate further, for better or worse.
  \item \textbf{Depth / hierarchical features:} MLPs can stack many hidden layers for hierarchical
  feature learning; RBF networks are inherently shallow.
\end{itemize}

\textbf{Applications.} RBF networks suit low-to-moderate-dimensional function
approximation/interpolation with clear local cluster structure, and settings needing fast, reliable
training with few hyperparameters. MLPs suit high-dimensional inputs (e.g.\ images), large and
complex datasets, and tasks that benefit from deep hierarchical feature extraction or good
extrapolation behavior.

\section*{Assignment 28}

\begin{itemize}
  \item \textbf{Completely Uninformed:} with little or no prior knowledge about the data, the
  input space is simply covered \emph{uniformly} with centers (e.g.\ a regular grid, or randomly
  placed with even density).
  \item \textbf{Input Data Driven:} centers are positioned/adjusted according to the
  \emph{distribution of the input data itself}, independent of the teacher values. Variants:
  heuristically picking a subset of the training points as centers, unsupervised clustering
  (e.g.\ $k$-means) to find data-driven centers, or self-organized approaches (e.g.\ a
  self-organizing map) producing the centers.
  \item \textbf{Error Driven:} centers (and widths) are adjusted with respect to the
  \emph{output error}, i.e.\ supervised learning -- e.g.\ gradient descent on the global error $F$,
  using exactly the $\partial F/\partial C_k$ and $\partial F/\partial s_k$ derived in
  assignment 32.
\end{itemize}

\section*{Assignment 29}

Fixed centers/widths, so only the $K{\times}M$ output weight matrix $\mathbf G$ (entries
$g_{k,m}$; the lecture calls the same object $\mathbf W$) must be found, from the $P{\times}K$
matrix $\mathbf R$ of RBF outputs ($R_{pk}={}^pr_k$) via
\[
  \mathbf R^{+} = (\mathbf R^{T}\mathbf R)^{-1}\mathbf R^{T}, \qquad \mathbf G = \mathbf R^{+}\hat{\mathbf Y}
\]
Counting multiplications for each step (assuming Gauss--Jordan elimination for the matrix
inverse, which needs $K^3$ multiplications for a $K{\times}K$ matrix):

\begin{enumerate}
  \item \textbf{Build $\mathbf R$} ($P{\times}K$ entries): each entry ${}^pr_k$ needs the squared
  Euclidean distance over $N$ dimensions ($N$ multiplications) plus one multiplication to divide
  by $2s_k^2$: $N+1$ multiplications per entry.
  \[
    P\cdot K\cdot(N+1)
  \]
  \item \textbf{$\mathbf R^{T}\mathbf R$} ($K{\times}K$ result, each entry a dot product of length
  $P$):
  \[
    K^2\cdot P
  \]
  \item \textbf{Invert the $K{\times}K$ matrix} (Gauss--Jordan):
  \[
    K^3
  \]
  \item \textbf{$\mathbf R^{+}=(\mathbf R^{T}\mathbf R)^{-1}\mathbf R^{T}$} ($K{\times}P$ result,
  each entry a dot product of length $K$):
  \[
    K\cdot P\cdot K = K^2\cdot P
  \]
  \item \textbf{$\mathbf G=\mathbf R^{+}\hat{\mathbf Y}$} ($K{\times}M$ result, each entry a dot
  product of length $P$):
  \[
    K\cdot M\cdot P
  \]
\end{enumerate}

Summing all five steps:
\[
  P K(N{+}1) + K^2P + K^3 + K^2P + KMP = PK(N{+}1{+}2K{+}M) + K^3
\]
\[
  \boxed{\text{total multiplications} = P\,K\,(N+2K+M+1) + K^3}
\]

\section*{Assignment 30}

Aspects to report when an AI system was used, extending the set from the previous sheets:
\begin{enumerate}
  \item Which AI system(s) were used (name, and model/version if known).
  \item Which prompt(s) were used (verbatim, or a close summary).
  \item What percentage of the final solution the AI contributed.
  \item Whether/why/when/how often the AI session had to be reset to get a usable answer.
  \item Which AI-generated answers were discarded, and why.
  \item \textbf{How did you check that the AI answer is correct?} (e.g.\ hand-recomputation,
  cross-checking against the lecture slides or the primary literature, running/testing code,
  consulting a second independent source.)
  \item \emph{(proposed additions)} Which parts of the AI's draft were edited or rewritten by hand
  afterwards, and why.
  \item Any limitations or open uncertainties that remain even after verification.
  \item Whether the AI was given the course material (slides/assignment text) directly, or
  answered purely from its own general training knowledge -- relevant for judging reliability.
\end{enumerate}

\section*{Assignment 31}

\textbf{AI use:} I used Claude (Anthropic) to work through this part, as the sheet requires.
Prompt: the assignment text, asking for a formula-based proposal to extend RBF neurons from radial
to elliptical (axis-dependent) characteristics. AI contribution: the formulas and explanation below
are as given by the AI; I verified them by checking that setting all $s_{nk}$ equal collapses the
elliptical formula back to the standard radial/Gaussian one exactly, and that the level sets
$z_k=\text{const}$ are indeed axis-aligned ellipses. No re-prompting or discarded answers were
needed. (Aspects reported per assignment 30, items 1-6 above.)

\textbf{What changes.} Standard RBF neurons use one shared width $s_k$ for the squared distance in
\emph{all} directions equally:
\[
  d_k^2 = \|{}^p\mathbf X - \mathbf C_k\|^2 = \sum_{n=1}^N ({}^px_n-c_{nk})^2, \qquad
  z_k = \exp\!\left(-\frac{d_k^2}{2s_k^2}\right)
\]
so level sets of $z_k$ are perfect (hyper-)spheres around $\mathbf C_k$. To get an elliptical
characteristic, give every RBF neuron $k$ \emph{one width per input axis},
$s_{1k},\dots,s_{Nk}$, instead of a single $s_k$, and replace the distance term by a per-axis
weighted sum:
\[
  \boxed{d_k^2 = \sum_{n=1}^N \frac{({}^px_n-c_{nk})^2}{s_{nk}^2}}, \qquad
  z_k = \exp\!\left(-\tfrac12\, d_k^2\right) = \exp\!\left(-\tfrac12\sum_{n=1}^N \frac{({}^px_n-c_{nk})^2}{s_{nk}^2}\right)
\]
Equivalently, in matrix form with $\mathbf S_k=\mathrm{diag}(s_{1k}^2,\dots,s_{Nk}^2)$:
\[
  z_k = \exp\!\left(-\tfrac12 ({}^p\mathbf X-\mathbf C_k)^T \mathbf S_k^{-1} ({}^p\mathbf X-\mathbf C_k)\right)
\]
Level sets $z_k=\text{const}$ are now axis-aligned ellipses/ellipsoids, with semi-axis $n$
proportional to $s_{nk}$: a direction with a larger $s_{nk}$ stretches the ellipse along that axis.
Setting all $s_{nk}=s_k$ recovers the original radial formula exactly, so this is a strict
generalization.

\textbf{What part of the RBF approach needs to change:} each RBF neuron now needs $N$ width
parameters instead of $1$ (i.e.\ $K\!\cdot\!N$ widths in total instead of $K$), and the distance
formula inside the transfer function changes as above. If widths are adjusted error-driven (gradient
descent, assignment 28/32), the single derivative $\partial F/\partial s_k$ is replaced by $N$
separate partial derivatives $\partial F/\partial s_{nk}$, one per axis. Nothing else about the
network (3-layer structure, linear output layer, centers $\mathbf C_k$) needs to change. (A further
generalization -- not required here -- replaces the diagonal $\mathbf S_k$ with a full
symmetric positive-definite matrix $\boldsymbol\Sigma_k$, allowing ellipses rotated at arbitrary
angles, not just axis-aligned ones.)

\section*{Assignment 32}

\[
  F = \frac12\sum_{p=1}^P\sum_{m=1}^M\left({}^p\hat y_m-{}^py_m\right)^2, \qquad
  z_k({}^p\mathbf X) = \exp\!\left(-\frac{\|{}^p\mathbf X-\mathbf C_k\|^2}{2s_k^2}\right), \qquad
  {}^py_m = \sum_{k=1}^K w_{km}\,z_k({}^p\mathbf X)
\]

Write ${}^p\delta_m:={}^p\hat y_m-{}^py_m$ (output layer is linear, so $f'=1$, no extra factor).
Since ${}^py_m$ depends on $z_k({}^p\mathbf X)$ only through $w_{km}$:
\[
  \frac{\partial F}{\partial z_k({}^p\mathbf X)} = \sum_{m=1}^M {}^p\delta_m\cdot(-w_{km}) = -\sum_{m=1}^M {}^p\delta_m\,w_{km}
\]

\textbf{Derivative w.r.t.\ $\mathbf C_k$.} With $\|{}^p\mathbf X-\mathbf C_k\|^2=\sum_n({}^px_n-c_{nk})^2$:
\[
  \frac{\partial z_k({}^p\mathbf X)}{\partial c_{nk}} = z_k({}^p\mathbf X)\cdot\frac{{}^px_n-c_{nk}}{s_k^2}
  \quad\Longrightarrow\quad
  \frac{\partial z_k({}^p\mathbf X)}{\partial \mathbf C_k} = z_k({}^p\mathbf X)\cdot\frac{{}^p\mathbf X-\mathbf C_k}{s_k^2}
\]
Chaining through all $P$ patterns:
\[
  \boxed{\frac{\partial F}{\partial \mathbf C_k} = -\sum_{p=1}^{P}\left(\sum_{m=1}^M {}^p\delta_m\,w_{km}\right) z_k({}^p\mathbf X)\,\frac{{}^p\mathbf X-\mathbf C_k}{s_k^2}}
\]

\textbf{Derivative w.r.t.\ $s_k$.}
\[
  \frac{\partial z_k({}^p\mathbf X)}{\partial s_k} = z_k({}^p\mathbf X)\cdot\frac{\|{}^p\mathbf X-\mathbf C_k\|^2}{s_k^3}
\]
\[
  \boxed{\frac{\partial F}{\partial s_k} = -\sum_{p=1}^{P}\left(\sum_{m=1}^M {}^p\delta_m\,w_{km}\right) z_k({}^p\mathbf X)\,\frac{\|{}^p\mathbf X-\mathbf C_k\|^2}{s_k^3}}
\]

Both gradients have the same structure as standard Backpropagation: a backpropagated error term
$\sum_m {}^p\delta_m w_{km}$ (analogous to a hidden-layer $\delta$), multiplied by the local
sensitivity of $z_k$ to $\mathbf C_k$ or $s_k$ -- here coming from differentiating the Gaussian
itself, instead of a transfer-function derivative $f'(\mathrm{net})$.

\end{document}
